\section{Preliminary}
\label{sec:preliminary}

To design an attribute-guided pruning framework, it is essential to understand how intrinsic Gaussian parameters relate to their perceptual and structural significance. Each Gaussian primitive $\mathcal{G}_i = \{x, y, z, \text{SH}_{0\text{-}47}, o, s_0, s_1, s_2, r_0, r_1, r_2, r_3\}$ encodes its spatial position, view-dependent color, opacity, scale, and rotation.  We observe that the relative importance of Gaussians can be largely inferred from their attributes and local neighborhood relations. 

Four observations reveal the primitive importance: (1) Existing rendering-based methods already provide evidence that opacity and volume are central to significance estimation: by weighting projected areas with these factors, they approximate each Gaussian’s contribution to blending and occlusion. Consistent with this, primitives with \textbf{small scales or low opacity} generally have limited visual influence compared to dominant neighbors, as shown in the top-right region of Fig.~\ref{fig:gs_redundancy}.
(2) Spatial isolation provides another cue for importance estimation. Gaussians are typically aligned with the underlying geometry, while those with \textbf{abnormally large distances to their $k$ nearest neighbors} tend to have lower importance, as illustrated by the floating points in the top-left region of Fig.~\ref{fig:gs_redundancy}.
(3) Under-optimization offers the third evidence for importance prediction. During densification, only Gaussians that remain consistently visible to the training cameras receive stable updates, while those with insufficient visibility fail to converge. Such Gaussians often exhibit an \textbf{abnormal appearance}: some display inconsistent or random colors compared with their neighbors, as seen in the clutter beneath the train (bottom-left and bottom-right of Fig.~\ref{fig:gs_redundancy}); others appear as nearly uniform single-color blobs because their higher-order spherical harmonics remain close to zero, leaving only the low-order terms to dominate the color. Both cases indicate insufficient optimization and, thus, low perceptual importance.
(4) Finally, both global and local magnitude statistics provide complementary signals for importance estimation—\textbf{local relative magnitudes capture occlusion and redundancy, while global magnitudes balance visibility across the scene}.

These observations motivate the construction of a compact feature representation that embeds intrinsic Gaussian attributes with local neighborhood statistics.
Building on this representation, we model the relationship between these features and Gaussian importance. Since explicitly hand-crafting this mapping is challenging, we adopt a lightweight MLP to learn it automatically, as described in Sec.~\ref{sec:method}.